\uuid{gIUn}
\titre{Matrice d'une application linéaire}
\auteur{hanani ; mzari}
\organisation{exo7}
\licence{CC BY-NC-SA 4.0}
\datecreate{2020-01-20}
\niveau{L1}
\module{Algèbre}
\chapitre{Matrice}
\sousChapitre{Matrice et application linéaire}
\difficulte{1}
\theme{Matrice d'une application linéaire}
\exo7qcmid{lille-2.matrice-application.002}

\contenu{
\begin{question}
\qclassification{108.03}
\qtags{motcle=Matrice d'une application linéaire}
On considère $\Rr$ et $\Rr^2$ munis de leurs bases canoniques et $f$ l'application linéaire définie par :
$$\begin{array}{rccc}f:&\Rr &\to& \Rr^2\\
& x&\to &(x,-x).  \end{array}$$
La matrice de $f$ relativement aux bases canoniques  est :
\begin{answers}
\bad{$\left(\begin{array}{rc}
1&-1\\
\end{array}\right).$}
\good{$\left(\begin{array}{r}
1\\
-1\\
\end{array}\right).$}
\bad{$\left(\begin{array}{rc}
1&-1\\
0&0\\
\end{array}\right).$}
\bad{$\left(\begin{array}{rcc}
1&0\\
-1&0
\end{array}\right).$}
\end{answers}
\begin{explanations} Soit ${\cal {B}} =\{1\}$ et ${\cal {B}}'=\{e_1,e_2\}$ les bases canoniques de $\Rr$ et $\Rr^2$ respectivement. Comme $f(1)= (1,-1)$, la matrice de $f$  relativement à ces bases est la matrice : $\left(\begin{array}{rc}1\\-1\end{array}\right).$
\end{explanations}
\end{question}
}
